\documentclass[12pt]{jlreq}
\usepackage{amsmath, amssymb}

\newcommand{\C}{\mathbb{C}}
\newcommand{\R}{\mathbb{R}}
\newcommand{\Z}{\mathbb{Z}}
\newcommand{\N}{\mathbb{N}}
\newcommand{\F}{\mathbb{F}}
\newcommand{\Prob}{\mathbb{P}}
\newcommand{\Exp}{\mathbb{E}}
\newcommand{\dd}{\,d}
\newcommand{\Ker}{\operatorname{Ker}}
\newcommand{\Impart}{\operatorname{Im}}
\newcommand{\Hom}{\operatorname{Hom}}

\begin{document}

\(\R\) 上の実数値関数列 \(f_n,\ n=1,2,\ldots,\) と実数値関数 \(f\) が与えられ，次の \(3\) 条件をみたすとする。

\begin{enumerate}
  \item[(a)] 各 \(f_n\) は単調非減少，すなわち \(x<y\) ならば \(f_n(x)\le f_n(y)\) である。
  \item[(b)] \(f\) は \(\R\) 上の連続関数である。
  \item[(c)] 任意の \(\varphi\in C_0(\R)\) に対して
  \[
    \lim_{n\to\infty}\int_{\R}f_n(x)\varphi(x)\dd x
    =\int_{\R}f(x)\varphi(x)\dd x
  \]
  が成立する。ただし \(C_0(\R)\) はコンパクトな台をもつ \(\R\) 上の連続関数全体を表す。
\end{enumerate}

このとき，\(n\to\infty\) とすれば，\(f_n\) は \(f\) に \(\R\) 上広義一様収束することを示せ。また，条件 (b), (c) を満たし結論が成立しないような関数列 \(f_n\) と関数 \(f\) の例をあげよ。

\end{document}